% !TeX program = xelatex
% !TEX root = main.tex
% !TeX encoding = UTF-8
\documentclass[10.5pt,compsoc]{CjC}
%\usepackage{CJKutf8}
%\usepackage{CJK}
\usepackage{graphicx}
\usepackage{footmisc}
\usepackage{subfigure}
\usepackage{url}
\usepackage{multirow}
\usepackage[noadjust]{cite}
\usepackage{amsmath,amsthm}
\usepackage{amssymb,amsfonts}
\usepackage{booktabs}
\usepackage{tabularx}
\usepackage{color}
\usepackage{ccaption}
\usepackage{float}
\usepackage{fancyhdr}
\usepackage{caption}
\usepackage{xcolor,stfloats}
\usepackage{comment}
\setcounter{page}{1}
\graphicspath{{figures/}}
\usepackage{cuted}
\usepackage{captionhack}
\usepackage{epstopdf}

%===============================%
% XeLaTeX + CJK font setup. CjC.cls is based on ctexart; these explicit fonts
% reduce the risk of garbled Chinese output on Windows TeX Live.
\setmainfont{Times New Roman}
\setCJKmainfont{SimSun}
\setCJKsansfont{SimHei}
\setCJKmonofont{FangSong}

% Manual body and heading spacing overrides for the compact CjC two-column
% layout. These settings intentionally override the class defaults so that long
% Chinese section headings never visually collide with adjacent paragraphs.
\setlength{\parskip}{0.38\baselineskip plus 0.08\baselineskip minus 0.04\baselineskip}
\makeatletter
\renewcommand\section{\@startsection {section}{1}{\z@}%
  {4.0ex \@plus 1.0ex \@minus .2ex}%
  {2.4ex \@plus .5ex}%
  {\normalfont\zihao{4-}\heiti\bfseries\raggedright}}
\renewcommand\subsection{\@startsection{subsection}{2}{\z@}%
  {3.2ex \@plus .8ex \@minus .2ex}%
  {1.8ex \@plus .4ex}%
  {\normalfont\zihao{5}\heiti\bfseries\raggedright}}
\renewcommand\subsubsection{\@startsection{subsubsection}{3}{\z@}%
  {2.6ex \@plus .6ex \@minus .2ex}%
  {1.4ex \@plus .3ex}%
  {\normalfont\zihao{5}\songti\bfseries\raggedright}}
\makeatother

\headevenname{\mbox{\quad} \hfill  \mbox{\zihao{-5}{\songti  计\quad \quad 算\quad \quad 机\quad \quad 学\quad \quad 报 } \hspace {50mm} \mbox{\songti  投稿稿件 }}}%
\headoddname{\songti  投稿稿件 \hfill
多模态大模型在智能电网中的应用综述 }%

\renewcommand{\thefootnote}{\fnsymbol{footnote}}
\setcounter{footnote}{0}
\renewcommand\footnotelayout{\zihao{5-}}

\newtheoremstyle{mystyle}{0pt}{0pt}{\normalfont}{1em}{\bf}{}{1em}{}
\theoremstyle{mystyle}
\renewcommand\figurename{figure~}
\renewcommand{\tablename}{表}
\captionsetup[table]{name=表,labelsep=quad}
\makeatletter
\renewcommand{\fnum@table}{表~\thetable}
\makeatother
\renewcommand{\thesubfigure}{(\alph{subfigure})}
\newcommand{\upcite}[1]{\textsuperscript{\cite{#1}}}
\renewcommand{\labelenumi}{(\arabic{enumi})}
\newcommand{\tabincell}[2]{\begin{tabular}{@{}#1@{}}#2\end{tabular}}
\newcommand{\abc}{\color{white}\vrule width 2pt}
\makeatletter
\renewcommand{\@biblabel}[1]{[#1]\hfill}
\makeatother
\setlength\parindent{2em}

\begin{document}
\hyphenpenalty=50000
\makeatletter
\newcommand\mysmall{\@setfontsize\mysmall{7}{9.5}}
\newenvironment{tablehere}
  {\def\@captype{table}}

\let\temp\footnote
\renewcommand \footnote[1]{\temp{\zihao{-5}#1}}

\thispagestyle{plain}%
\thispagestyle{empty}%
\pagestyle{CjCheadings}

\begin{table*}[!t]
\vspace {-13mm}
\begin{tabular}{p{168mm}}
\zihao{5-}
\songti
投稿稿件 \hfill 计\quad 算\quad 机\quad 学\quad 报\hfill Submission Manuscript
\zihao{5-}\songti
投稿阶段卷期页码由编辑部填写 \hfill CHINESE JOURNAL OF COMPUTERS \hfill To be assigned \\
\hline\\[-4.5mm]
\hline\end{tabular}

\centering
\vspace{11mm}
\heiti
{\zihao{2} 多模态大模型在智能电网中的应用综述}

\vskip 5mm

{\zihao{3} \fangsong
作者信息投稿时补充$^{1)}$
}

\vspace{5mm}
\zihao{6}{\songti
$^{1)}$(作者单位、通信作者与基金信息投稿时补充)
}

\zihao{6}{\heiti
本文为投稿前匿名/半匿名稿件，作者、单位、基金及联系方式信息按投稿系统要求补充。
}

\vskip 5mm
{\centering
\begin{tabular}{p{160mm}}
\zihao{5-}{
\setlength{\baselineskip}{16pt}\selectfont{
\noindent\heiti 摘\quad 要\quad   \songti
智能电网在源网荷储协同、设备状态感知、调度运行和安全治理中产生图像、视频、红外、文本、时序、拓扑和知识图谱等多源异构数据，传统单模态智能方法难以支撑统一理解、交互推理和可信决策。面向这一问题，本文综述多模态大模型在智能电网中的应用进展，并提出``数据模态--模型机制--电力任务--可信约束''四维分类框架，统一刻画输入来源、能力实现、业务目标和部署边界之间的关系。在此基础上，本文梳理视觉语言模型、指令微调、检索增强生成和知识图谱等关键技术，总结电力巡检、故障诊断、调度辅助、应急演练和安全监管等典型应用，分析知识增强可信推理、数据治理、物理约束、实时性、可解释性和安全责任等挑战，并展望电力专用多模态基础模型、物理约束融合、动态知识增强、人在回路决策和标准化评测等方向。
\par}}\\[2mm]

\zihao{5-}{\noindent
\heiti 关键词  \quad \songti  多模态大模型；智能电网；视觉语言模型；电力巡检；故障诊断；调度辅助决策；知识增强
}\\[2mm]
\zihao{5-}{\heiti 中图法分类号 	\songti
TP391 \rm{\quad \quad \quad     }
\heiti DOI号: \songti
投稿时不提供DOI号}
\end{tabular}}

\vskip 7mm

\begin{center}
\zihao{3}{ {\heiti A Survey on Applications of Multimodal Large Models in Smart Grids}}\\
\vspace {5mm}
\zihao{5}{ {\heiti Course Paper Author$^{1)}$}}\\
\vspace {2mm}
\zihao{6}{\heiti {$^{1)}$(Course Paper, affiliation omitted)}}

\end{center}

\begin{tabular}{p{160mm}}
\zihao{5}{
\setlength{\baselineskip}{18pt}\selectfont{
{\bf Abstract}\quad
Smart grids generate heterogeneous data from source-grid-load-storage coordination, equipment-state perception, dispatching operation, and safety governance, including images, videos, infrared sensing, text, time series, topology, and knowledge graphs. Conventional single-modal artificial intelligence methods are insufficient for unified understanding, interactive reasoning, and trustworthy decision support. This paper surveys multimodal large models for smart grids and proposes a four-dimensional taxonomy of data modality, model mechanism, power task, and trustworthy constraint to characterize the relationships among input sources, capability realization, business objectives, and deployment boundaries. Based on this framework, the paper reviews key techniques such as vision-language models, instruction tuning, retrieval-augmented generation, and knowledge graphs; summarizes representative applications in power inspection, fault diagnosis, dispatching assistance, emergency rehearsal, and safety supervision; analyzes challenges in knowledge-enhanced trustworthy reasoning, data governance, physical constraints, real-time requirements, explainability, and safety responsibility; and discusses future directions including power-domain multimodal foundation models, physical-constraint integration, dynamic knowledge enhancement, human-in-the-loop decision making, and standardized evaluation.
\par}}\\

\setlength{\baselineskip}{18pt}\selectfont{
\zihao{5}{\noindent
\vspace {5mm}
{\bf Keywords}\quad multimodal large models; smart grid; vision-language models; power inspection; fault diagnosis; dispatching assistance; knowledge enhancement
}\par}
\end{tabular}

\setlength{\tabcolsep}{2pt}
\begin{tabular}{p{0.05cm}p{16.15cm}}
\multicolumn{2}{l}{\rule[4mm]{40mm}{0.1mm}}\\[-3mm]
&\begin{songti}
投稿前稿件：作者、单位、基金、邮箱及手机号等投稿专用信息按期刊要求补充。
\end{songti}
\end{tabular}\end{table*}
\clearpage\clearpage
\begin{strip}
\vspace {-13mm}
\end{strip}
\heiti
\zihao{5}
\setlength{\baselineskip}{17pt}\selectfont
\vskip 3mm


\section{引言}

\songti
新型电力系统以高比例新能源接入、源网荷储协同、分布式资源互动和数字化调度为重要特征。作为其数字化、智能化实现形态，智能电网不仅需要处理潮流、负荷、状态量等结构化数据，还需要融合巡检图像、红外图像、作业视频、设备文本、拓扑图结构和运行规程等多源异构信息。已有研究表明，大语言模型正在进入电力系统规划、运行、市场、运维和网络安全等环节\cite{ghafari2025comprehensive,mirshekali2025review,niu2025llm_power}。然而，智能电网中的智能化任务并非单纯的文本生成、问答或单一模态识别问题，其关键在于将视觉缺陷、设备语义、拓扑位置、运行状态和规程约束关联到同一决策链条中，并在多模态感知、领域知识检索、状态推理和安全校核之间形成可追溯的业务闭环。

多模态大模型（Multimodal Large Language Model, MLLM）为上述问题提供了新的研究路径。以 Transformer 架构\cite{vaswani2017attention}为基础，大语言模型通过规模化预训练形成语言理解、上下文推理和指令交互能力\cite{brown2020language}；视觉语言模型进一步通过图文对齐与视觉指令微调，使图像内容能够与自然语言描述、问答和解释过程相连接\cite{radford2021learning,liu2023visual}。在此基础上，MLLM 支持多模态输入、自然语言解释和跨任务泛化，有助于将感知识别、知识检索和交互推理组织为统一建模过程。在电力领域，视觉语言模型和多模态指令微调方法为巡检图像识别、缺陷文本生成、规程知识查询、故障处置辅助和调度交互提供了统一建模的可能\cite{li2024foundation_power,xiang2025mllm_power}。

然而，多模态大模型不能被简单理解为通用模型在电力数据上的直接迁移。电力系统具有高安全性要求、强实时性约束、显式物理约束和清晰责任边界等特点，模型输出需要满足设备约束、运行规程、调度安全边界和可审计要求。现有综述已从电力系统大语言模型、能源系统大模型和电力大模型等角度对相关研究进行了梳理\cite{sarwar2025large,fan2025cognitive,zhang2026large,du2026key}，为理解大模型在电力领域的应用奠定了基础。然而，已有综述多从语言模型能力、单一应用场景或能源系统总体趋势出发，重点讨论模型能力、应用场景和发展挑战。对于智能电网中的多模态任务，仍缺少一个能够同时连接数据来源、模型机制、业务任务和可信约束的分析框架。具体而言，现有研究尚未系统回答以下问题：不同电力模态如何被 MLLM 表征与对齐；多模态感知、检索增强和物理约束推理如何共同服务于巡检、诊断、调度和运维等任务；模型输出在进入电力生产系统前应满足哪些安全、可靠和可审计要求。

为此，本文围绕多模态大模型在智能电网中的应用展开综述，提出``数据模态--模型机制--电力任务--可信约束''四维系统分析框架。其中，数据模态回答模型处理什么信息，模型机制回答多模态信息如何被对齐、检索和推理，电力任务回答模型服务于哪些业务环节，可信约束回答模型输出如何满足安全、物理和工程边界。本文主要贡献包括：
（1）构建``数据模态--模型机制--电力任务--可信约束''四维分析框架，用于统一刻画 MLLM 在智能电网中的应用路径；
（2）建立典型场景与模型能力之间的映射关系，分析 MLLM 在感知、诊断、调度、知识增强和人机协同中的作用边界；
（3）从可靠性、安全性、可解释性、评测体系和工程部署等方面归纳 MLLM 进入电力生产系统前必须满足的可信要求。

\begin{table*}[t]
\centering
\caption{现有大模型与电力系统相关综述的比较}
\label{tab:survey_comparison}
\zihao{6}\songti
\setlength{\tabcolsep}{3pt}
\renewcommand{\arraystretch}{1.18}
\begin{tabularx}{\textwidth}{p{0.12\textwidth} p{0.17\textwidth} c c c c X}
\toprule
文献 & 关注对象 & 多模态数据归纳 & 模型机制分析 & 电网任务映射 & 可信约束讨论 & 主要局限 \\
\midrule
Ghafari等\cite{ghafari2025comprehensive} & 电力系统大语言模型应用 & 较少 & 部分 & 系统 & 部分 & 侧重 LLM 在电力系统不同业务环节中的应用梳理，对图像、红外、拓扑等多模态数据及其统一建模关系讨论较少。 \\
Mirshekali等\cite{mirshekali2025review} & 能源系统大语言模型应用、挑战与前景 & 较少 & 部分 & 部分 & 部分 & 关注能源系统总体应用与发展挑战，智能电网多模态数据来源和生产级任务约束不是其分析主线。 \\
Sarwar等\cite{sarwar2025large} & 电力系统 LLM 应用文献综述 & 较少 & 部分 & 系统 & 部分 & 较系统覆盖电力系统 LLM 应用，但主要围绕语言模型能力展开，对多模态表征、对齐和物理约束推理的关联分析不足。 \\
Fan等\cite{fan2025cognitive} & 面向认知智能的电力系统 LLM & 部分 & 部分 & 部分 & 部分 & 强调 LLM 支撑电力系统认知智能的研究脉络，但尚未系统展开多模态数据、模型机制和可信边界之间的对应关系。 \\
Zhang等\cite{zhang2026large} & 能源系统中的大语言模型 & 部分 & 系统 & 部分 & 部分 & 对能源系统 LLM、增强技术和未来方向归纳较完整，但研究对象偏能源系统总体趋势，智能电网多模态任务映射仍不充分。 \\
杜博骏等\cite{du2026key} & 电力系统大模型关键科学问题与挑战 & 部分 & 部分 & 部分 & 系统 & 侧重电力系统大模型的科学问题、挑战和发展方向，对具体多模态数据类型与典型任务之间的映射仍需细化。 \\
向思屿等\cite{xiang2025mllm_power} & 电力领域多模态大模型应用 & 系统 & 部分 & 部分 & 部分 & 已关注电力多模态大模型应用，但对模型机制、业务任务和生产系统可信约束的统一框架化分析仍有扩展空间。 \\
本文 & 多模态大模型在智能电网中的应用 & 系统 & 系统 & 系统 & 系统 & 侧重综述与框架构建，具体算法实现和大规模实证评测仍有待后续研究。 \\
\bottomrule
\end{tabularx}
\end{table*}

本文余下部分组织如下：第2节介绍多模态大模型关键技术基础；第3节构建智能电网多模态数据与任务体系；第4节总结典型应用场景；第5节讨论知识增强、RAG与物理约束推理；第6节分析工程部署、评测体系与安全可信挑战；第7节展望未来研究方向；第8节给出结论。


\section{多模态大模型关键技术基础}

多模态大模型并不是单一技术的简单组合，而是由语言建模、跨模态对齐、指令微调、知识增强、工具调用和领域适配等机制共同构成的模型体系。语言模型提供语义理解与指令交互底座，视觉语言模型完成图文语义对齐，多模态指令微调将基础能力转化为任务响应，知识增强与工具调用提供外部依据和校验能力，电力领域适配决定模型能否进入电力生产系统。对于智能电网而言，这些技术的意义不只在于提升图像识别或文本问答能力，更在于将巡检图像、红外测温、作业视频、时序量测、拓扑关系、设备语义和规程知识组织到统一的证据链和推理链中。因此，本节从语言模型、视觉语言模型、多模态指令微调、知识增强与工具调用、电力领域适配五个方面梳理 MLLM 的关键技术基础，为后文分析智能电网中的多模态任务、知识增强可信推理和工程部署约束奠定基础。

\subsection{大语言模型：从语言理解到指令交互}

Transformer 架构通过自注意力机制实现了对长序列依赖关系的建模，成为当前大语言模型和多模态大模型的基础结构\cite{vaswani2017attention}。在大规模语料预训练的基础上，大语言模型表现出上下文学习、少样本迁移、指令跟随和自然语言生成能力\cite{brown2020language}。这些能力使模型能够在不针对每一类任务单独设计特征工程的情况下，完成规程问答、工单摘要、报告生成和故障处置建议等知识密集型任务\cite{ghafari2025comprehensive,niu2025llm_power}。

然而，通用语言模型本质上仍以文本序列为主要输入，其能力边界取决于训练语料、上下文窗口和生成机制。在智能电网中，许多关键任务并不以纯文本形式存在，而是由图像、时序量测、拓扑结构、设备台账和运行规程共同构成。例如，故障诊断不仅需要理解现场描述，还需要结合保护动作记录、设备状态量和历史检修记录；巡检任务也不仅是识别图像中的缺陷，还需要将缺陷位置、设备类别和规程要求关联起来。因此，大语言模型为电力知识理解和人机交互提供了基础，但难以单独支撑智能电网中的跨模态感知与可信决策。

\subsection{视觉语言模型：从图文对齐到多模态感知}

视觉语言模型通过图像与文本之间的语义对齐，使视觉内容能够被自然语言描述、检索和解释。CLIP 利用自然语言监督学习可迁移视觉表示，主要提升跨模态检索和零样本视觉识别能力\cite{radford2021learning}。在此基础上，LLaVA 等代表性视觉语言模型通过连接视觉编码器与大语言模型，并利用视觉指令微调增强模型的图像问答、描述生成和交互式解释能力\cite{liu2023visual}。二者分别代表了从表征对齐到视觉指令交互的技术演进，也构成电力视觉语言模型的主要技术来源。

对于智能电网而言，视觉语言模型的价值在于打通``视觉证据''和``业务语义''之间的联系。电力巡检图像中的绝缘子破损、导线异物、表计异常和设备外观缺陷，并不是孤立的视觉类别，而是与设备名称、部件层级、缺陷等级、巡检规程和处置建议相关联的业务对象。视觉语言模型能够将图像中的设备状态、缺陷特征和部件信息与自然语言语义进行关联，并进一步与规程、工单和设备台账建立联系，从而为缺陷识别、异常解释和巡检报告生成提供统一的建模基础。PowerVLM 等电力视觉语言模型研究也表明，视觉语言模型正在向电力专用场景适配方向发展\cite{ouyang2026powervlm}。

需要指出的是，当前主流视觉语言模型主要围绕图像--文本模态展开，而智能电网还包含红外图像、作业视频、时序量测、拓扑图结构和知识图谱等模态。红外和视频数据需要关注时间连续性与异常区域定位，时序量测需要保留动态变化和运行边界，拓扑结构则需要表达节点、支路和设备连接关系。因而，电力 MLLM 不能简单等同于图文 VLM，而需要结合专用编码器、时序建模、图结构表示或外部工具调用来扩展模态边界\cite{bernier2025powergraph,zhengxiang2025data,du2026key}。

\subsection{多模态指令微调：从感知识别到交互推理}

多模态指令微调是将通用视觉语言能力转化为任务交互能力的重要环节。与传统监督学习直接学习固定标签不同，指令微调通过构造图像--文本--指令--响应数据，使模型学习根据自然语言要求完成识别、描述、问答、解释和生成等多种任务。经过多模态指令微调后，模型不再只是输出类别标签，而可以围绕同一输入证据生成缺陷描述、回答追问、解释判断依据，并给出后续处理建议\cite{liu2023visual,xiang2025mllm_power}。

这种能力对于智能电网尤其重要。电力业务场景通常不是一次性分类任务，而是包含“观察--解释--追问--确认--处置”的交互流程。例如，在巡检场景中，运维人员可能不仅关心图像中是否存在缺陷，还会进一步询问缺陷属于哪类设备、是否影响运行安全、是否需要立即处理以及依据哪条规程。多模态指令微调使 MLLM 能够以自然语言接口连接视觉证据和业务问题，从而支持更接近真实工作流程的人机协同分析。

不过，多模态指令微调也存在明显局限。首先，电力领域高质量图文指令数据稀缺，缺陷样本通常存在长尾分布和标注不一致问题。其次，模型生成的解释可能具有语言流畅性，但不一定具备物理正确性和规程一致性。再次，通用多模态数据中的视觉语义与电力设备语义存在差异，直接迁移可能导致设备部件识别错误或缺陷等级判断偏差。因此，面向智能电网的多模态指令微调需要结合领域知识、业务流程和安全约束进行专门设计。

\subsection{知识增强与工具调用：从生成回答到可验证推理}

生成式模型容易受到参数记忆不足、知识过期和幻觉问题影响，因此知识增强是 MLLM 进入电力领域的重要技术路径。检索增强生成（retrieval-augmented generation, RAG）通过将外部文档检索结果作为上下文输入，使模型能够基于规程、标准、设备台账、历史工单和应急预案生成回答\cite{lewis2020retrieval}。与单纯依赖模型参数相比，RAG 更适合电力系统中具有明确版本、责任和依据要求的知识密集型任务。

知识图谱则提供了另一种结构化知识组织方式。它能够显式表达设备、部件、缺陷、工况、规程和处置动作之间的关系，为多模态大模型提供可查询、可追溯和可推理的领域知识空间\cite{ji2022survey}。在智能电网场景中，知识图谱可用于连接图像中的设备部件、台账中的设备编码、规程中的专业术语和工单中的故障描述，从而缓解不同数据源之间语义不一致的问题。层次化图检索增强生成等研究也表明，图结构知识能够进一步增强复杂知识检索和生成过程\cite{shen2026hgrag}。

除 RAG 和知识图谱外，工具调用也是提高模型可靠性的重要机制。对于调度辅助、故障处置和运行状态分析等任务，模型不能只根据语言推理生成建议，而应调用潮流计算、规则校验、仿真平台、优化器或专家系统进行验证\cite{zhangjun2023operation,yang2025operator,li2024risks}。在该机制中，MLLM 主要承担语义理解、证据组织、候选方案生成和结果解释，潮流计算、规则校验和优化求解则作为外部可信工具提供可验证约束。由此，MLLM 的角色应从``直接给出控制指令''转变为``组织证据、生成候选方案、调用工具校验并解释结果''的辅助决策模块。ReAct 和 AutoGen 等智能体框架也表明，推理、行动和工具调用可以被组织为可交互的任务求解流程\cite{yao2023react,wu2023autogen}。对于高风险电力任务，这种``生成--检索--校验--复核''的链式机制比单次生成回答更符合工程部署要求。

\subsection{电力领域适配：从通用 MLLM 到可信电力智能系统}

通用 MLLM 不能直接等同于可部署的电力智能系统。电力系统具有高安全性要求、强实时性约束、显式物理约束和清晰责任边界，模型输出需要满足设备限值、调度规程、安全校核和可审计要求。因此，电力领域适配不仅包括数据和模型层面的微调，还包括数据治理、知识增强、系统部署和安全管控等层面的联合设计。

从数据角度看，电力多模态数据存在隐私敏感、标注成本高、跨场景差异大和共享受限等问题。巡检图像可能包含设备铭牌、站内位置和作业人员信息，调度日志可能包含运行方式和事故处置过程，直接集中训练或外部云端推理可能带来安全风险\cite{li2024risks}。联邦学习、私有化部署、数据脱敏和权限控制可以在一定程度上缓解数据安全与模型能力之间的矛盾\cite{ouyang2026powervlm}，但也会引入模型更新、性能评估和系统运维的新问题。与此同时，电力数据质量、模态对齐和跨场景一致性仍会直接影响模型适配效果\cite{zhengxiang2025data}。

从模型角度看，电力领域适配需要同时考虑参数高效微调、模型压缩、边云协同和持续更新。边缘巡检终端、变电站现场设备和调度辅助系统对时延、算力和网络稳定性有不同要求，因此模型规模、推理方式和部署位置需要结合业务风险分级进行选择。电力智算基础设施、边云协同和行业模型部署能力将直接影响 MLLM 的工程可用性\cite{fangfang2025computing,du2026key}。对于巡检报告生成、规程问答和应急演练等低风险或半在线场景，可以优先采用 RAG 和人机协同部署；对于调度辅助和故障处置等高风险任务，则应引入物理模型、仿真平台和人工确认机制。此时，MLLM 更适合作为证据组织、候选建议生成和结果解释模块，而不应直接触发控制动作。
总体来看，MLLM 在智能电网中的关键价值不只是融合多种输入模态，而是将感知证据、领域知识、推理过程和可信约束连接到同一业务链条中。后续章节将基于“数据模态--模型机制--电力任务--可信约束”四维框架，进一步分析智能电网中的典型任务、应用场景和可信部署挑战。

\section{智能电网多模态数据与任务体系}


\subsection{多模态数据类型}

智能电网的多模态数据可分为七类：一是可见光图像和视频，用于设备外观、线路通道和作业现场识别；二是红外和局部放电等状态感知数据，用于反映设备热异常或绝缘缺陷；三是电压、电流、功率、频率等时序量测数据；四是电网拓扑、设备连接和站线关系等图结构数据；五是规程、标准、调度指令、检修票和工单等文本数据；六是设备台账、缺陷库和检修记录等业务数据；七是由上述数据抽象形成的电力知识图谱。不同模态在采集频率、噪声水平、语义粒度和安全等级上存在差异，决定了多模态融合不能仅靠简单拼接，而需要结合任务目标和业务约束进行建模\cite{li2024foundation_power,du2026key,zhengxiang2025data}。

表\ref{tab:data_modality}从数据来源、业务语义、适配机制和可信约束四个方面归纳上述模态，突出``数据模态''维度并为后续任务映射提供输入侧依据。

\begin{table*}[t]
\centering
\caption{智能电网多模态数据类型与建模难点}
\label{tab:data_modality}
\zihao{6}\songti
\setlength{\tabcolsep}{3pt}
\renewcommand{\arraystretch}{1.16}
\begin{tabularx}{\textwidth}{p{0.13\textwidth} p{0.20\textwidth} p{0.21\textwidth} p{0.20\textwidth} X}
\toprule
数据模态 & 典型来源 & 业务语义 & 适配机制 & 建模难点与可信约束 \\
\midrule
可见光图像与视频 & 无人机、巡检机器人、固定摄像头、作业现场视频 & 设备外观、线路通道、缺陷区域、作业行为 & 视觉编码器、图文对齐、视觉指令微调 & 光照、遮挡、视角和天气变化显著，需控制跨场站泛化、误报漏报和人员隐私风险 \\
红外与局部放电等状态感知数据 & 红外测温、局部放电检测、声纹与振动传感器 & 热异常、绝缘劣化、机械状态和早期故障征兆 & 专用传感编码器、异常区域定位、跨模态对齐 & 数据噪声和环境温度影响较大，需结合设备负荷、运行工况和规程阈值解释异常 \\
时序量测数据 & 电压、电流、功率、频率、保护动作记录、告警序列 & 运行状态、扰动过程、负荷变化和故障演化 & 时序建模、状态估计、检索增强和工具校验 & 采样频率不一、缺失与异常值较多，模型输出需满足实时性、物理一致性和安全边界 \\
拓扑与图结构数据 & 电网接线图、站线关系、设备连接、馈线结构 & 节点--支路关系、设备位置、供电路径和拓扑约束 & 图表示学习、图增强语言模型、知识图谱 & 拓扑随运行方式变化，需保持设备实体对齐、结构更新和潮流约束可验证 \\
文本与规程数据 & 运行规程、技术标准、调度指令、检修票、缺陷描述 & 业务规则、操作流程、处置依据和责任要求 & LLM、RAG、指令微调、引用生成 & 知识版本和适用范围需明确，回答必须可追溯、可审计并避免引用过期条款 \\
业务台账与历史记录 & 设备台账、缺陷库、检修记录、历史工单、事故案例 & 设备身份、生命周期状态、历史缺陷和处置经验 & 实体对齐、案例检索、知识图谱增强 & 数据格式异构且质量不均，需解决编码标准、权限控制和跨系统一致性问题 \\
电力知识图谱 & 由设备、部件、缺陷、工况、规程和处置动作抽象形成 & 领域实体、关系约束、可解释推理路径 & 图检索、图推理、RAG 和规则约束 & 构建和更新成本高，实体关系错误会放大到检索、推理和生成环节 \\
\bottomrule
\end{tabularx}
\end{table*}


\subsection{多模态任务谱系}

从任务层次看，智能电网多模态大模型可支持感知识别、语义理解、知识检索、状态推理、方案生成和人机协同六类能力。感知识别主要面向巡检图像缺陷检测和作业现场识别；语义理解关注图像、文本和设备台账之间的关系解释；知识检索用于规程问答和证据定位；状态推理面向故障诊断、运行态势研判和异常原因分析；方案生成用于抢修方案、巡检报告和应急演练脚本；人机协同则强调模型建议经人工复核后进入业务闭环\cite{ghafari2025comprehensive,mirshekali2025review,xiang2025mllm_power}。


\subsection{面向智能电网的分类框架}

本文建议从``数据模态--模型机制--电力任务--可信约束''四个维度理解多模态大模型应用。数据模态刻画输入来源，模型机制说明模型如何完成跨模态对齐、知识检索、推理生成和工具调用，电力任务对应巡检、诊断、调度、应急和运维等业务目标，可信约束则限定模型输出能否进入生产流程。四个维度不是简单并列关系，而是构成从``输入证据''到``模型能力''、再到``业务动作''和``部署边界''的分析链条：同一模型机制在不同数据模态和电力任务下可能产生不同风险，同一业务任务也必须根据实时性、物理可行性、责任归属和可审计性设置不同的使用边界。

图\ref{fig:framework}展示了多模态大模型赋能智能电网的五层技术架构，用于说明从数据层、模型层、知识层、应用层到可信治理层的工程实现链路。与之不同，本文在第3.3节提出的``数据模态--模型机制--电力任务--可信约束''四维分析框架，主要用于比较不同应用场景在输入来源、能力实现、业务目标和部署边界上的差异。该四维框架贯穿全文后续结构：第4节按照电力任务维度比较巡检、故障诊断、调度辅助和应急演练等典型应用；第5节重点展开模型机制维度中的 RAG、知识图谱和可信推理；第6节从可信约束维度讨论数据治理、物理约束、可解释性、评测和系统集成问题；第7节则围绕四个维度归纳电力专用模型、物理约束融合、动态知识增强、人在回路和标准化评测等研究方向。该框架有助于避免将多模态大模型泛化为单纯的图像识别或文本问答技术，也有助于区分不同场景的成熟度与风险等级。

表\ref{tab:model_mechanisms}进一步从``模型机制''维度总结 MLLM 的关键技术及其电力适配方式，强调基础模型能力需要通过知识增强、工具校验和领域部署机制才能进入智能电网业务链条。

\begin{table*}[t]
\centering
\caption{MLLM 关键技术机制及其电力适配}
\label{tab:model_mechanisms}
\zihao{6}\songti
\setlength{\tabcolsep}{3pt}
\renewcommand{\arraystretch}{1.16}
\begin{tabularx}{\textwidth}{p{0.15\textwidth} p{0.17\textwidth} p{0.22\textwidth} p{0.24\textwidth} X}
\toprule
技术机制 & 机制定位 & 支撑的电力模态/任务 & 电力适配方式 & 主要约束 \\
\midrule
语言建模与指令交互 & 语义理解底座 & 规程、工单、台账、调度指令与运行报告 & 通过领域语料、提示模板和指令数据增强规程问答、报告生成和故障处置建议能力\cite{brown2020language,ghafari2025comprehensive} & 难以直接处理非文本证据，且可能受幻觉、上下文窗口和知识过期影响 \\
图文与跨模态对齐 & 感知证据接入 & 巡检图像、红外图像、视频帧与设备文本 & 通过视觉语言模型将缺陷区域、设备部件和业务语义对齐\cite{radford2021learning,liu2023visual,ouyang2026powervlm} & 以图像--文本为主，对时序和拓扑数据需额外编码或工具支持 \\
多模态指令微调 & 任务响应转换 & 缺陷识别、图像问答、报告草拟和交互式解释 & 构造图像--文本--指令--响应样本，使模型适应``观察--解释--追问--处置''流程\cite{xiang2025mllm_power} & 依赖高质量电力标注和指令数据，需控制缺陷等级、规程依据和解释一致性 \\
RAG 与知识库增强 & 外部依据接入 & 规程问答、应急预案、历史工单和标准条款检索 & 将规程、标准、台账和预案作为外部证据输入生成过程\cite{lewis2020retrieval,liwenchao2025rag} & 检索证据需满足版本有效、来源权威、引用可追溯和冲突可识别 \\
知识图谱与图增强推理 & 结构化关系建模 & 设备--部件--缺陷关系、电网拓扑、工况与处置动作 & 用图谱或图增强语言模型表达实体关系、拓扑位置和业务约束\cite{ji2022survey,bernier2025powergraph,shen2026hgrag} & 实体对齐和图谱更新成本高，结构错误会影响下游检索与推理 \\
工具调用与物理校验 & 可验证推理闭环 & 调度辅助、故障处置、运行状态分析和方案校核 & 调用潮流计算、规则校验、仿真平台、优化器或专家系统验证候选建议\cite{yao2023react,yang2025operator,zhangjun2023operation} & MLLM 应承担证据组织和解释功能，不能越过外部校验直接触发控制动作 \\
领域部署与持续适配 & 生产系统落地 & 边缘巡检、站端分析、调度辅助和知识服务 & 结合参数高效微调、模型压缩、边云协同、私有化部署和权限审计\cite{ouyang2026powervlm,fangfang2025computing,du2026key} & 受数据安全、算力资源、响应时延、模型更新和系统运维约束 \\
\bottomrule
\end{tabularx}
\end{table*}

\begin{figure*}[htbp]
\centering
\fbox{\begin{minipage}{0.92\textwidth}
\centering
\begin{tabular}{c}
\fbox{数据层：图像、视频、红外、文本、时序、拓扑、知识图谱}\\[1mm]
$\Downarrow$\\[-1mm]
\fbox{模型层：视觉语言对齐、指令微调、跨模态表示、工具调用}\\[1mm]
$\Downarrow$\\[-1mm]
\fbox{知识层：规程库、设备台账、历史工单、RAG、知识图谱}\\[1mm]
$\Downarrow$\\[-1mm]
\fbox{应用层：巡检识别、故障诊断、调度辅助、应急演练}\\[1mm]
$\Downarrow$\\[-1mm]
\fbox{可信治理层：物理约束、安全校验、证据追溯、人在回路}
\end{tabular}
\end{minipage}}
\caption{多模态大模型赋能智能电网的五层技术架构。图中五层描述从数据到可信治理的工程实现链路；第3.3节提出的四维分析框架用于理解不同应用场景在数据模态、模型机制、电力任务和可信约束上的差异。}
\label{fig:framework}
\end{figure*}


\subsection{与传统电力人工智能方法的关系}

多模态大模型并不替代传统机器学习、优化算法和电力机理模型，而是对其形成补充。传统方法在潮流计算、状态估计、故障分类和优化调度中具有明确数学结构和可验证性；多模态大模型则更适合开放语义理解、复杂文档交互、多源证据融合和业务流程协同。对于高安全等级的调度辅助与运行控制类任务，应将模型输出定位为辅助建议，并通过物理模型、仿真平台和调度规则进行约束校验\cite{zhangjun2023operation,li2024risks}。


\section{多模态大模型在智能电网中的典型应用}

本节按照第3.3节提出的四维框架，从输入模态、模型机制、任务输出和可信约束四个方面分析典型应用场景。需要强调的是，以下应用中的多模态大模型主要承担感知理解、证据组织、候选建议和结果解释等辅助功能，其输出仍需接受业务规则、物理模型和人工复核的约束。

\subsection{电力巡检与设备缺陷识别}

电力巡检是多模态大模型具有较好落地基础的场景。传统视觉模型通常聚焦于设备或缺陷的检测分类，而多模态大模型可以将图像理解、缺陷语义解释和巡检结果描述连接起来。例如，电力视觉语言模型可根据巡检图像识别绝缘子破损、导线异物、表计异常和设备外观缺陷，并生成面向巡检任务的文本描述\cite{ouyang2026powervlm,wang2025inspection,liu2026defect}。在巡检机器人场景中，自然语言交互还可用于任务说明和结果解释，但任务执行仍应受巡检路径、作业许可和现场安全规则约束。

该方向的主要瓶颈在于跨场站泛化、复杂天气鲁棒性、缺陷样本不均衡和误报漏报责任边界。由于电力巡检图像往往涉及不同设备厂家、拍摄角度和现场环境，单一数据集上的性能难以代表真实部署效果。因此，后续研究需要建设覆盖多设备、多场景和多模态证据的评测集，并将模型输出与人工复核、工单系统和缺陷闭环治理结合。


\subsection{故障诊断、设备运维与抢修方案生成}

故障诊断和设备运维要求模型从多源证据中推断故障类型、故障位置和处置策略。与单纯分类任务不同，电力故障处置通常包含现场现象描述、保护动作记录、设备状态量、历史工单、图像证据和规程条款。大模型指令微调可用于生成抢修方案和运维建议\cite{deng2026repair}，类 DeepSeek 范式等推理模型也被用于探索电力装备故障诊断中的推理链组织\cite{jiang2025deepseek}。已有故障诊断综述指出，机器学习和大语言模型为新型电力系统故障诊断提供了新的方法补充，但仍需与机理知识和专家经验结合\cite{zhangjianliang2025fault}。

从工程角度看，故障诊断类应用尤其需要证据链。模型不应只给出``可能是某故障''的结论，还应明确依据哪些量测、图像、规程或历史案例，并给出不确定性提示。对于抢修方案生成，应将模型输出限制为辅助建议，经现场人员和调度人员审核后执行，避免生成式模型直接触发高风险操作。


\subsection{调度辅助、运行状态感知与辅助决策}

调度辅助是电力大模型应用中价值高但风险也最高的方向之一。现有研究开始评估大语言模型在电力系统调度任务中的能力\cite{mongaillard2024scheduling}，也有工作探索将大语言模型作为配电网电压控制中的操作员角色\cite{yang2025operator}。PowerGraph-LLM 等图增强方法表明，将电网拓扑结构和语言模型结合有助于提升运行决策中的结构化理解能力\cite{bernier2025powergraph}。中文研究也从运行控制、推理大模型和配电系统专用大模型角度讨论了智能决策前景\cite{zhangjun2023operation,yutao2026reasoning}。

但是，调度业务具有强实时、强约束和强安全属性，通用模型不能绕过潮流约束、设备限值和调度规程。更合理的路径是将多模态大模型用于调度语义理解、态势摘要、异常提示、预案检索、候选方案说明和校验结果解释，再由优化算法、仿真系统和人工调度员完成最终校验。因此，当前阶段应强调辅助决策和方案解释，而非完全自主控制。


\subsection{应急演练、安全监管与工程实践}

应急演练和安全监管属于知识密集、流程密集但相对可控的应用场景。基于高级 RAG 与多模态大模型的电力应急演练研究表明，模型可将应急预案、现场图像、处置流程和问答交互结合起来，生成演练脚本、风险提示和处置建议\cite{liwenchao2025rag}。在电力工程实践中，行业大模型已被用于智能巡检、安全督察、调度辅助和知识服务等场景\cite{ghafari2025comprehensive,li2024foundation_power}。

这类场景的优势在于可先以离线或半在线方式部署，降低直接控制风险。其主要挑战是知识库版本管理、业务流程一致性和输出责任归属。若模型引用过期规程或错误工单，可能造成误导。因此，应急与监管类应用也需要建立证据追溯、版本控制和人工确认机制。

表\ref{tab:task_mapping}按照四维框架对上述典型应用进行归纳，其中``MLLM 能力定位''用于说明模型在业务流程中可承担的辅助环节，``可信约束与部署边界''用于限定其进入生产系统的条件。表\ref{tab:representative_studies}进一步归纳已有代表性研究，突出不同工作在数据模态、模型机制和电力任务上的侧重点。

\begin{table*}[htbp]
\centering
\caption{典型电力任务与 MLLM 能力映射}
\label{tab:task_mapping}
\zihao{6}\songti
\setlength{\tabcolsep}{3pt}
\renewcommand{\arraystretch}{1.16}
\begin{tabularx}{\textwidth}{p{0.14\textwidth} p{0.20\textwidth} p{0.22\textwidth} p{0.21\textwidth} X}
\toprule
电力任务 & 主要数据模态 & MLLM 能力定位 & 典型输出 & 可信约束与部署边界 \\
\midrule
电力巡检与缺陷识别 & 可见光图像、视频、红外、设备台账、巡检规程 & 图像理解、缺陷语义解释、跨模态检索和报告草拟 & 缺陷候选、图像证据说明、缺陷等级建议、巡检报告初稿 & 需人工复核缺陷等级和处置建议，重点控制跨场景泛化、误报漏报和隐私泄露风险 \\
故障诊断与设备运维 & 状态量、保护记录、告警、历史工单、图像和规程文本 & 多源证据汇聚、诊断依据组织、案例检索和方案生成 & 故障类型候选、证据链、抢修方案草案、备品备件提示 & 应定位为辅助诊断，需标注不确定性并经现场人员、专家系统或调度人员确认 \\
调度辅助与运行态势感知 & 时序量测、拓扑结构、告警序列、调度指令、运行规程 & 态势摘要、预案检索、候选方案解释、工具调用和校验结果解释 & 异常提示、运行状态摘要、预案条款、候选调度方案说明 & 不应直接进入控制闭环，必须经过潮流计算、设备限值、规则校验和人工调度员确认 \\
规程问答与知识服务 & 规程标准、设备台账、检修票、缺陷库、知识图谱 & RAG、知识图谱增强、引用生成和语义检索 & 条款定位、操作流程解释、知识问答、证据来源列表 & 需保证知识库版本有效、引用可追溯、权限可控，并识别证据冲突和适用范围 \\
应急演练与安全监管 & 应急预案、现场图像、流程文本、角色信息、历史案例 & 场景生成、流程推演、风险提示和人机协同问答 & 演练脚本、处置步骤、风险清单、监管记录摘要 & 适合离线或半在线部署，需维持流程一致性、责任归属和人工确认机制 \\
规划、市场与综合决策辅助 & 负荷预测、市场文本、政策文件、拓扑和设备约束 & 文档理解、方案对比、约束摘要和解释性生成 & 规划方案摘要、市场规则解释、风险因素清单 & 受数据时效、政策边界和模型不确定性影响，输出应作为分析材料而非自动决策结果 \\
\bottomrule
\end{tabularx}
\end{table*}

\begin{table*}[htbp]
\centering
\caption{代表性应用研究归纳}
\label{tab:representative_studies}
\zihao{6}\songti
\setlength{\tabcolsep}{3pt}
\renewcommand{\arraystretch}{1.14}
\begin{tabularx}{\textwidth}{p{0.15\textwidth} p{0.13\textwidth} p{0.18\textwidth} p{0.18\textwidth} p{0.18\textwidth} X}
\toprule
文献 & 应用场景 & 数据模态 & 模型机制 & 主要贡献 & 相对于本文框架的局限 \\
\midrule
欧阳旭东等\cite{ouyang2026powervlm} & 电力视觉语言模型与巡检 & 巡检图像、文本描述 & VLM、联邦学习、模型剪枝 & 探索电力专用视觉语言模型及轻量化适配路径 & 主要聚焦图像--文本场景，对调度、拓扑和生产级可信约束覆盖有限 \\
王启其等\cite{wang2025inspection} & 电力巡检图像识别 & 巡检图像、缺陷文本 & 多模态大模型识别与描述 & 面向巡检图像缺陷识别提供应用示例 & 任务边界集中在巡检识别，对跨任务机制和可信部署讨论不足 \\
刘爽\cite{liu2026defect} & 设备缺陷识别 & 图像、多模态特征 & 多模态融合、数据增强 & 关注复杂场景下设备缺陷识别与数据增强 & 偏向识别模型与数据增强，尚未形成 MLLM 决策链条分析 \\
邓养吾等\cite{deng2026repair} & 故障抢修方案生成 & 故障文本、工单和业务知识 & 指令微调、方案生成 & 将大模型用于故障抢修方案智能生成 & 生成结果仍需证据追溯、规程校验和人工确认 \\
江秀臣等\cite{jiang2025deepseek} & 电力装备故障诊断 & 设备状态、故障知识和文本证据 & 推理模型、诊断推理链 & 探索类 DeepSeek 范式在装备诊断中的应用模式 & 主要强调推理范式，对图像、拓扑和多模态对齐机制涉及较少 \\
Mongaillard等\cite{mongaillard2024scheduling} & 电力系统调度 & 调度任务文本、运行约束描述 & LLM 调度能力评估 & 评估 LLM 在调度任务中的能力边界 & 以语言模型评估为主，尚需接入物理仿真和多模态运行证据 \\
Yang等\cite{yang2025operator} & 配电网电压控制 & 配电网状态、动作空间和控制反馈 & LLM 操作员、工具/环境交互 & 探索 LLM 参与电压控制任务的可行性 & 高风险控制场景需更严格的物理校验、权限隔离和人在回路机制 \\
Bernier等\cite{bernier2025powergraph} & 电力系统图增强推理 & 电网拓扑、节点支路关系和文本 & 图增强大语言模型 & 将图结构信息引入电力系统语言模型推理 & 对视觉、红外和作业视频等感知模态覆盖不足 \\
李文超等\cite{liwenchao2025rag} & 电力应急演练 & 应急预案、现场图像、流程文本 & RAG、多模态问答 & 将高级 RAG 与 MLLM 用于应急演练场景 & 适合流程和知识服务，仍需强化证据版本、流程一致性和责任边界 \\
\bottomrule
\end{tabularx}
\end{table*}


\section{知识增强、RAG与电力领域可信推理}


\subsection{电力知识图谱与多模态知识组织}

电力知识图谱能够将设备、部件、缺陷、规程、状态量、检修记录和运行场景组织为结构化知识网络。与纯文本知识库相比，知识图谱能够显式表达实体关系和业务约束，有利于模型进行可解释推理\cite{ji2022survey}。在变电站智能运检、电力领域多模态知识图谱和能源数字网络知识图谱等研究中，知识图谱被用于支撑设备状态理解、跨模态检索和知识服务\cite{fuxuhui2025mmkg,zhouaihua2026kg}。

多模态知识图谱的关键难点在于实体对齐和动态更新。图像中的设备部件、台账中的设备编码、规程中的专业术语和工单中的故障描述可能对应同一实体或状态。如果对齐错误，后续检索和推理都会受到影响。因此，电力知识图谱建设需要结合标准化命名、设备编码体系、图像标注和业务流程规则。


\subsection{RAG在电力规程和业务知识中的应用}

RAG 将外部知识检索与模型生成结合，适合电力规程问答、故障处置建议和应急演练等知识密集型任务\cite{lewis2020retrieval,liwenchao2025rag}。在电力场景中，RAG 的评价重点不应只是答案是否流畅，还应关注检索证据是否权威、版本是否有效、引用是否可追溯以及不同证据之间是否冲突。层次化图检索等方法可进一步提高复杂知识结构下的检索组织能力\cite{shen2026hgrag}。


\subsection{可信推理机制}

可信推理要求模型输出从``生成文本''转向``给出可验证建议''。在故障诊断中，模型应说明使用了哪些图像证据、状态量、历史工单和规程条款；在调度辅助中，模型应说明建议是否经过潮流约束、设备限值和安全规则校验；在应急演练中，模型应给出处置流程与预案来源。图增强语言模型、知识图谱和 RAG 可以共同构成证据链，为多模态大模型提供可审计的推理基础\cite{bernier2025powergraph,ji2022survey,li2024risks}。


\section{工程部署、安全可信与评测挑战}


\subsection{数据治理与隐私安全}

电力数据具有关键基础设施属性，涉及企业生产、用户隐私和电网安全。多模态大模型训练和部署面临数据孤岛、标注成本高、数据质量不一致和共享限制等问题\cite{zhengxiang2025data,du2026key}。例如，巡检图像可能包含设备铭牌、站内位置和作业人员信息，调度日志可能包含运行方式和事故处置信息，若直接进入集中式训练或外部云端推理，将带来敏感数据泄露风险。联邦学习、私有化部署、模型压缩和权限控制可在一定程度上缓解数据安全与模型能力之间的矛盾\cite{ouyang2026powervlm}，但也会带来模型更新、性能评估和系统运维的新问题。


\subsection{可靠性、实时性与物理约束}

智能电网中的模型输出必须符合电力系统物理约束。多模态大模型可能生成语义上合理但物理上不可行的建议，例如在配电网电压越限场景中建议切换无功补偿装置，却忽略相邻节点电压、线路潮流和设备投切次数约束；在故障处置中建议转供负荷，却未校验备用线路热稳定限值。智能电网大模型风险研究指出，模型在电网环境中的误用可能导致安全、隐私和运行风险\cite{li2024risks}。因此，面向调度辅助和故障处置的系统应建立``模型建议--规则校验--仿真验证--人工确认''的闭环，而不能将生成结果直接转化为控制指令。


\subsection{可解释性、可追溯性与责任边界}

可解释性是多模态大模型进入电力生产流程的前提。巡检识别需要说明缺陷位置和图像证据，故障诊断需要说明推理依据，调度辅助需要说明约束校验结果。若模型输出无法追溯来源，现场人员和调度人员难以判断其可信度。例如，模型若给出``疑似套管过热''结论，应同时给出红外图像区域、温升依据、历史缺陷记录和规程条款来源，而不能只输出自然语言判断。RAG、知识图谱和日志审计可在一定程度上提高输出可追溯性，但责任边界仍需制度化设计，尤其是当模型参与抢修方案生成或运行建议时\cite{lewis2020retrieval,deng2026repair,li2024risks}。


\subsection{评测体系与基准数据集缺失}

当前电力多模态大模型缺少统一基准测试，不同研究在数据来源、任务定义、评价指标和场景复杂度上差异较大。单纯使用准确率或召回率难以评价模型的业务价值。例如，巡检模型即使在单一数据集上具有较高缺陷识别准确率，也可能在强光、雨雾、遮挡或不同厂家设备场景下出现性能下降；调度问答模型即使回答流畅，也可能引用过期规程或忽略安全校验。面向智能电网的评测应同时覆盖识别准确性、跨场景鲁棒性、证据可追溯性、响应时延、安全违规率、幻觉率和人工采纳率。对于高风险任务，还应引入仿真环境和红队测试，评估模型在异常输入、冲突证据和规程变更情况下的表现\cite{ghafari2025comprehensive,mirshekali2025review,li2024risks}。


\subsection{工程部署与系统集成}

多模态大模型落地不仅是模型问题，也是系统工程问题。电力场景需要考虑边缘端算力、网络带宽、模型更新、业务系统接口、审计日志和安全运维流程。例如，边缘巡检终端可能无法承载大规模视觉语言模型的实时推理，应采用模型压缩、边云协同或离线批处理；调度辅助系统则需要与能量管理系统、仿真平台和权限审计模块隔离集成。电力智算系统和行业大模型实践表明，算力基础设施、数据治理平台和业务流程重构是模型规模化应用的基础\cite{fangfang2025computing,du2026key}。对于巡检和应急等场景，可优先采用离线分析或人机协同方式部署；对于调度辅助等高风险场景，应在仿真和沙箱环境中逐步验证。

表\ref{tab:challenge}从可信约束维度归纳上述挑战，将数据治理、多模态对齐、证据追溯、物理校核、实时部署、责任边界和评测体系分别对应到可操作的治理机制与未来研究方向。

\begin{table*}[htbp]
\centering
\caption{可信、安全与工程部署挑战矩阵}
\label{tab:challenge}
\zihao{6}\songti
\setlength{\tabcolsep}{3pt}
\renewcommand{\arraystretch}{1.16}
\begin{tabularx}{\textwidth}{p{0.14\textwidth} p{0.23\textwidth} p{0.19\textwidth} p{0.23\textwidth} X}
\toprule
可信约束维度 & 主要风险 & 影响任务/模态 & 建议校验与治理机制 & 后续研究方向 \\
\midrule
数据治理与隐私安全 & 数据孤岛、敏感信息泄露、标注不足、跨场景数据质量不一致 & 巡检图像、调度日志、工单台账、用户侧数据 & 数据分级分类、脱敏、权限控制、联邦学习、私有化部署和数据质量评估 & 建设可共享的脱敏样本、行业数据规范和跨场景质量评价体系 \\
多模态对齐与数据质量 & 图像、红外、时序、拓扑和文本之间时间戳、实体编码和语义粒度不一致 & 巡检、诊断、拓扑推理、知识问答 & 统一设备编码、实体对齐、时间同步、模态质量评分和异常数据标记 & 发展面向电力设备和拓扑关系的跨模态基准与自动对齐方法 \\
幻觉控制与证据追溯 & 模型生成无依据结论、引用过期规程或无法解释诊断依据 & 规程问答、故障诊断、应急演练、报告生成 & RAG、知识图谱、证据链记录、引用版本管理和冲突证据检测 & 建立面向电力知识的可追溯生成评价指标和证据一致性检测方法 \\
物理一致性与安全校核 & 候选建议违反潮流约束、设备限值、操作规程或安全稳定边界 & 调度辅助、故障处置、运行控制、抢修方案 & 规则校验、潮流计算、仿真验证、优化器复核和人工调度员确认 & 探索 MLLM 与机理模型、数字孪生和安全约束优化的协同推理 \\
实时性与工程部署 & 模型规模、推理时延、网络带宽、边缘算力和系统接口限制 & 在线监测、边缘巡检、站端分析、调度辅助 & 模型压缩、边云协同、缓存检索、灰度发布、沙箱验证和日志审计 & 构建面向电力生产系统的轻量化 MLLM 和可观测部署框架 \\
可解释性与责任边界 & 输出责任不清、越权建议、人工复核缺位和审计链不完整 & 调度、抢修、安全监管、人机协同 & 人在回路、角色权限控制、操作留痕、责任分级和审批流程 & 形成模型建议、人工确认、系统执行之间的责任划分和审计标准 \\
评测体系与安全测试 & 缺少统一数据集、任务定义、鲁棒性测试和安全红队评估 & 全部多模态任务，尤其是高风险调度和故障处置 & 场景化基准、幻觉率、安全违规率、时延、人工采纳率和红队测试 & 建立覆盖数据模态、模型机制、电力任务和可信约束的综合评测体系 \\
\bottomrule
\end{tabularx}
\end{table*}


\section{未来研究方向}


\subsection{面向电力系统的专用多模态基础模型}

未来研究需要从通用 MLLM 调用走向电力专用多模态基础模型建设。该方向包括电力设备视觉语义库、缺陷图文对齐数据、时序--拓扑联合表示、电力规程指令集和跨场景评测集。专用模型既要提升设备与业务语义理解能力，也要避免对私有数据和高风险控制流程的不当暴露\cite{du2026key,li2024foundation_power,ouyang2026powervlm}。


\subsection{物理约束与机理模型融合}

智能电网应用不能只依赖统计相关性和语言推理。未来应探索多模态大模型与潮流计算、优化调度、数字孪生和仿真平台的融合，使模型输出在生成阶段或后处理阶段接受物理约束校验。图增强语言模型和配电网电压控制研究为这一方向提供了初步启示\cite{bernier2025powergraph,yang2025operator,zhangjun2023operation}。


\subsection{动态知识增强与持续学习}

电力规程、设备状态和运行方式会持续变化，静态知识库难以长期保持有效。未来需要动态 RAG、可更新知识图谱和持续学习机制，使模型能够处理新设备、新故障和新业务流程。同时，应建立知识版本管理和变更审计机制，确保模型引用的规程和证据处于有效状态\cite{lewis2020retrieval,ji2022survey,shen2026hgrag}。


\subsection{多智能体协同与人在回路决策}

智能体和多智能体系统是未来趋势，而不宜作为当前智能电网 MLLM 应用的主线。ReAct 将推理和行动结合起来\cite{yao2023react}，AutoGen 等多智能体框架展示了多角色协同完成复杂任务的可能性\cite{wu2023autogen}。在电力场景中，多智能体可模拟调度员、运维人员、安监人员和知识检索模块之间的协作，用于配网故障研判和应急会商\cite{cui2026multiagent,dou2026autonomous}。但其应用必须受权限控制、审计追踪、安全隔离和人在回路确认约束。


\subsection{可信、安全与标准化评测}

多模态大模型能否进入生产级智能电网，取决于可信评测和安全标准。未来应建立覆盖缺陷识别、跨模态检索、规程问答、诊断推理和调度辅助的公开或行业级基准，评价指标应包括准确性、鲁棒性、时延、可解释性、幻觉率、安全违规率和人工采纳率。对于高风险业务，还应引入红队测试和故障注入验证，识别模型在异常输入和冲突证据下的失效模式\cite{li2024risks,ghafari2025comprehensive,mirshekali2025review}。


\section{结论}

多模态大模型为智能电网处理多源异构数据、理解电力业务语义和支撑人机协同决策提供了新的技术路径。本文提出的四维分类框架表明，相关应用需要同时考察数据模态、模型机制、电力任务和可信约束。与传统单模态模型相比，MLLM 在巡检图像理解、故障证据汇聚、规程知识问答、应急演练和调度辅助中具有更强的跨模态表达和交互能力。然而，智能电网的强安全、强实时和强物理约束属性决定了 MLLM 不能以通用生成模型形态直接进入控制闭环。未来研究应围绕电力专用多模态基础模型、知识增强可信推理、物理约束融合、标准化评测、隐私保护和人在回路机制持续推进。总体而言，多模态大模型在智能电网中的价值并非替代电力机理模型和专家体系，而是通过跨模态理解与知识增强机制，提升复杂业务场景下的感知、解释、推理和协同能力。

\vspace {5mm}
\centerline
{\zihao{5}
\heiti 参~考~文~献}

\begingroup
\renewcommand{\refname}{}
\zihao{5-}
\bibliographystyle{unsrt}
\bibliography{references_clean}
\endgroup

\end{document}
